% problem-000354 8 青蒿素全合成
\section*{8. 青蒿素全合成}
⼀。我国著名有机合成化学家、中国科学院院⼠周维善教授在⻘蒿素的全合成⽅⾯做出了开创性的⼯作，他领导的研究⼩组于1983年完成了⻘蒿素的⾸次全合成。该研究组所采⽤的合成路线如下：

（图：）

\subsection*{8-1}
⻘蒿素分⼦中共有⼏个⼿性碳原⼦？写出内酯环上C-2、C-3、C-4和C-5的构型（⽤�/�表⽰）。

\subsection*{8-2}
⽤系统命名法命名起始原料⾹茅醛和中间体A（⽤�/�表⽰其构型）。

\subsection*{8-3}
画出中间体B、F和G的⽴体结构式（⽤楔线式表⽰）。

\subsection*{8-4}
中间体C在⾼位阻强碱LDA存在下与3-三甲基硅基-3-丁烯-2-酮进⾏Michael加成反应时，除了得到中间体D之外，还产⽣了⼀种副产物，它是D的⽴体异构体。试画出这种副产物的⽴体结构式（⽤楔线式表⽰）。

\subsection*{8-5}
写出由中间体D到中间体E转化的机理（⽤反应式表⽰）。
